% GENERATED FILE. Edit canonical lessons, then run npm run generate:books.
% canonical-source-hash: fnv1a64:198414d55355703f
% canonical-lessons: GE-C343-konsequenz, GE-C343-resultat, GE-C343-folglich, GE-C343-sodass, GE-C343-daher

\chapter{Die Konsequenzen}
\label{ch:ge-die-konsequenzen}
\hlchaptermodality{343}

\begin{chapteropening}
\textbf{By the end of this chapter:} \emph{I can talk about consequences and results, and link a cause to its result with folglich, sodass and daher.}

Complete the last lesson of chapter 343: I can talk about consequences and results, and link a cause to its result with folglich, sodass and daher.
\end{chapteropening}

\section[die Konsequenz]{die Konsequenz — the consequence; consistency}

\label{lesson:GE-C343-konsequenz}



\begin{quote}
Before the new one: say the German for to trigger, then the German for the trigger.
\end{quote}

\subsection*{You'll want to know: die Konsequenz}
\textbf{die Konsequenz} — \textquotedblleft{}the consequence; consistency\textquotedblright{}.

\begin{quote}
Wer zu spät kommt, muss die Konsequenzen tragen.
\end{quote}

\begin{quote}
\emph{Anyone who comes late has to bear the consequences.}
\end{quote}

Plural: \emph{die Konsequenzen}. \emph{Konsequenzen ziehen} — to draw conclusions and act on them.

\begin{cousinweb}
\emph{Konsequenz} came through French \emph{conséquence} from Latin \emph{cōnsequentia} \textquotedblleft{}a consequence, a natural succession\textquotedblright{} (EtymDB; Lewis \& Short), from \emph{cōnsequī} \textquotedblleft{}to follow after\textquotedblright{}.
\end{cousinweb}

\begin{grammarlens}[title={What you've built}]
One more word for telling, explaining and describing at B1.
\end{grammarlens}

\subsection*{Your turn}
\begin{itemize}
\raggedright
  \item \emph{Say it:} \emph{die Konsequenz}
  \item \emph{Say it:} \emph{die Konsequenz}, once more
  \item \emph{Say it:} a sentence of your own with \emph{die Konsequenz}
  \item \emph{Recall:} say the German for to trigger, then the German for the trigger, then say \emph{die Konsequenz} again
\end{itemize}

\begin{tcolorbox}[breakable,colback=teal!4,colframe=teal!35!black,title={Before you move on}]
What does die Konsequenz mean? (\textquotedblleft{}The consequence; consistency\textquotedblright{}.) Say it once more.
\end{tcolorbox}
\section[das Resultat]{das Resultat — the result}

\label{lesson:GE-C343-resultat}



\begin{quote}
Before the new one: say the German for the trigger, then the German for the consequence.
\end{quote}

\subsection*{You'll want to know: das Resultat}
\textbf{das Resultat} — \textquotedblleft{}the result\textquotedblright{}.

\begin{quote}
Das Resultat des Tests war sehr gut.
\end{quote}

\begin{quote}
\emph{The result of the test was very good.}
\end{quote}

Plural: \emph{die Resultate}. A little more formal than \emph{das Ergebnis}.

\begin{cousinweb}
\emph{Resultat} was borrowed from French \emph{résultat} (EtymDB), from Medieval Latin \emph{resultātum}, built on \emph{resultāre} \textquotedblleft{}to spring back, to rebound\textquotedblright{} (Lewis \& Short): what springs back from an action.
\end{cousinweb}

\begin{grammarlens}[title={What you've built}]
One more word for telling, explaining and describing at B1.
\end{grammarlens}

\subsection*{Your turn}
\begin{itemize}
\raggedright
  \item \emph{Say it:} \emph{das Resultat}
  \item \emph{Say it:} \emph{das Resultat}, once more
  \item \emph{Say it:} a sentence of your own with \emph{das Resultat}
  \item \emph{Recall:} say the German for the trigger, then the German for the consequence, then say \emph{das Resultat} again
\end{itemize}

\begin{tcolorbox}[breakable,colback=teal!4,colframe=teal!35!black,title={Before you move on}]
What does das Resultat mean? (\textquotedblleft{}The result\textquotedblright{}.) Say it once more.
\end{tcolorbox}
\section[folglich]{folglich — consequently, therefore}

\label{lesson:GE-C343-folglich}



\begin{quote}
Before the new one: say the German for the consequence, then the German for the result.
\end{quote}

\subsection*{You'll want to know: folglich}
\textbf{folglich} — \textquotedblleft{}consequently, therefore\textquotedblright{}.

\begin{quote}
Es gab keinen Strom. Folglich konnten wir nicht kochen.
\end{quote}

\begin{quote}
\emph{There was no electricity. Consequently, we could not cook.}
\end{quote}

The verb comes right after it: \emph{folglich konnten wir …}. Rather formal; in speech people often say \emph{deshalb}.

\begin{cousinweb}
\emph{folglich} is built with \emph{-lich} on the family of \emph{folgen} \textquotedblleft{}to follow\textquotedblright{}, inherited from Old High German \emph{folgēn} and Proto-Germanic \emph{fulgijaną} (EtymDB): what follows from something.
\end{cousinweb}

\begin{grammarlens}[title={What you've built}]
One more word for telling, explaining and describing at B1.
\end{grammarlens}

\subsection*{Your turn}
\begin{itemize}
\raggedright
  \item \emph{Say it:} \emph{folglich}
  \item \emph{Say it:} \emph{folglich}, once more
  \item \emph{Say it:} a sentence of your own with \emph{folglich}
  \item \emph{Recall:} say the German for the consequence, then the German for the result, then say \emph{folglich} again
\end{itemize}

\begin{tcolorbox}[breakable,colback=teal!4,colframe=teal!35!black,title={Before you move on}]
What does folglich mean? (\textquotedblleft{}Consequently, therefore\textquotedblright{}.) Say it once more.
\end{tcolorbox}
\section[sodass]{sodass — so that (as a result)}

\label{lesson:GE-C343-sodass}



\begin{quote}
Before the new one: say the German for the result, then the German for consequently.
\end{quote}

\subsection*{You'll want to know: sodass}
\textbf{sodass} — \textquotedblleft{}so that (as a result)\textquotedblright{}.

\begin{quote}
Es hat stark geschneit, sodass die Schule geschlossen war.
\end{quote}

\begin{quote}
\emph{It snowed heavily, so that the school was closed.}
\end{quote}

\emph{Sodass} sends the verb to the end: \emph{…, sodass die Schule geschlossen war}. Also written \emph{so dass}. It names a result, not an aim.

\begin{cousinweb}
\emph{sodass} is the joined spelling of \emph{so dass}, from \emph{so} and \emph{dass} \textquotedblleft{}that\textquotedblright{}; both spellings are correct today. EtymDB lists the word without a history, so this is its plain make-up.
\end{cousinweb}

\begin{grammarlens}[title={What you've built}]
One more word for telling, explaining and describing at B1.
\end{grammarlens}

\subsection*{Your turn}
\begin{itemize}
\raggedright
  \item \emph{Say it:} \emph{sodass}
  \item \emph{Say it:} \emph{sodass}, once more
  \item \emph{Say it:} a sentence of your own with \emph{sodass}
  \item \emph{Recall:} say the German for the result, then the German for consequently, then say \emph{sodass} again
\end{itemize}

\begin{tcolorbox}[breakable,colback=teal!4,colframe=teal!35!black,title={Before you move on}]
What does sodass mean? (\textquotedblleft{}So that (as a result)\textquotedblright{}.) Say it once more.
\end{tcolorbox}
\section[daher]{daher — therefore, that is why; from there}

\label{lesson:GE-C343-daher}



\begin{quote}
Before the new one: say the German for consequently, then the German for so that.
\end{quote}

\subsection*{You'll want to know: daher}
\textbf{daher} — \textquotedblleft{}therefore, that is why; from there\textquotedblright{}.

\begin{quote}
Ich hatte kein Geld dabei, daher habe ich nichts gekauft.
\end{quote}

\begin{quote}
\emph{I had no money on me, so I did not buy anything.}
\end{quote}

The verb comes right after it: \emph{daher habe ich …}. \emph{Daher kommt das!} — so that is where it comes from.

\begin{cousinweb}
\emph{daher} joins \emph{da} \textquotedblleft{}there\textquotedblright{} and \emph{her} \textquotedblleft{}hither, from\textquotedblright{}, and first meant \textquotedblleft{}from there\textquotedblright{}; the sense \textquotedblleft{}for that reason\textquotedblright{} grew out of it. EtymDB has the word but no history for it.
\end{cousinweb}

\begin{grammarlens}[title={What you've built}]
One more word for telling, explaining and describing at B1.
\end{grammarlens}

\subsection*{Your turn}
\begin{itemize}
\raggedright
  \item \emph{Say it:} \emph{daher}
  \item \emph{Say it:} \emph{daher}, once more
  \item \emph{Say it:} a sentence of your own with \emph{daher}
  \item \emph{Recall:} say the German for consequently, then the German for so that, then say \emph{daher} again
\end{itemize}

\begin{tcolorbox}[breakable,colback=teal!4,colframe=teal!35!black,title={Before you move on}]
What does daher mean? (\textquotedblleft{}Therefore, that is why; from there\textquotedblright{}.) Say it once more.
\end{tcolorbox}
